% Copyright (c) 2026 Marcio Luiz Lima de Oliveira and V. el Hakim.
% Distribution and modification are subject to LPPL version 1.3c.
% SPDX-License-Identifier: LPPL-1.3c
% Part of the LPPL-covered work listed in MANIFEST.txt.
% See LICENSE and LICENSING.txt for the licence and scope.

\chapter{References and your publications}\label{chap:05}
\renewcommand{\FigLocation}{Chapters/Chapter_5/fig/}
\renewcommand{\ChapLocation}{Chapters/Chapter_5/}
\begin{absbox}
Maintain two databases with different purposes. References record the literature consulted; publications record the author’s own output. A keyword determines which printed list receives an entry.
\end{absbox}

\pubfootnote{This chapter illustrates a publication-origin note. The accompanying
entry \cite{morganExample} is fictional and demonstrates formatting only.}

\section{Two files, two roles}
Put external literature in \texttt{References/references.bib} and your own
publications in \texttt{References/publications.bib}. Both files are loaded by
\texttt{main.tex}. Every entry intended for the Publications list needs
\texttt{keywords = \{mypub\}}. The database filename alone does not select
the printed list.
\begin{lstlisting}[language={}]
@article{morganExample,
  author       = {Morgan, Alex and Taylor, Jordan},
  title        = {An Example Publication for This Template},
  journaltitle = {Example Journal},
  year         = {2026},
  volume       = {1},
  pages        = {1--8},
  keywords     = {mypub},
  note         = {Fictional example; replace before submission}
}
\end{lstlisting}
The real documentation references are separated from these fictional entries.
For example, the \texttt{memoir} manual~\cite{memoirManual} is consulted
literature, while \cite{morganExample} is an illustrative author publication.
Do not duplicate the same entry key across the two databases.

\section{Printing the lists}
The end of \texttt{main.tex} uses the original keyword distinction:
\begin{lstlisting}
\printbibliography[notkeyword=mypub,heading=bibintoc]
\printbibliography[keyword=mypub,title={Publications},
  heading=bibintoc]
\end{lstlisting}
\texttt{heading=bibintoc} adds each heading to the table of contents. The
alphabetic citation labels are shared across the document; the split changes
membership, not the identity of a cited item. The default style is
\texttt{alphabetic}, not an IEEE bibliography style.

Only cited entries normally print. To include a particular publication that
is not cited in the prose, use an explicit key:
\begin{lstlisting}
\nocite{morganSupplement}
\end{lstlisting}
\nocite{morganSupplement}
Avoid \verb|\nocite{*}| unless you really want every entry from both databases.
An own-publication entry without \texttt{mypub} will appear in the literature
list; an external reference with that keyword will appear under Publications.

\section{Citations and publication notes}
Use \verb|\cite{key}| for a normal citation, \verb|\textcite{key}| when the author
belongs in the sentence, and \verb|\parencite{key}| for the style's parenthetical
form. An optional locator can be useful, for example
\verb|\cite[Chapter 3]{memoirManual}|. Citation commands inside the literal
listings do not count as citations.

The existing \verb|\pubfootnote{text}| helper inserts an unnumbered footnote
without advancing the ordinary footnote counter. Place it directly after a
chapter heading or its opening abstract when giving a paper's origin:
\begin{lstlisting}
\pubfootnote{Parts of this chapter are based on
  \cite{morganExample}.}
\end{lstlisting}
Use accurate wording for an adapted or reproduced paper, include the
coauthors where appropriate, and retain any required publisher credit. The
note is not a substitute for a bibliography entry. Ordinary notes use
\verb|\myfootnote{text}| or \verb|\footnote{text}|.

\section{Changing the citation style deliberately}
The \texttt{biblatex} options are in \texttt{Configuration.tex}. For numeric
citations change \texttt{style=alphabetic} to \texttt{style=numeric}; keep the
keyword filters. The \texttt{url}, \texttt{doi} and \texttt{isbn} options control
those fields in bibliography entries, not the identifiers on the title pages.
\texttt{backref=true} adds citing-page information. The sample online references
include their access links in \texttt{howpublished} because URLs are disabled
by the inherited bibliography options.

If you switch to \texttt{backend=biber}, use Biber rather than BibTeX and rebuild
from clean auxiliary files. Change the tool and the option together. The
\texttt{biblatex} manual~\cite{biblatexManual} describes the richer styles and
data model. Keep a single bibliography system; do not add \texttt{natbib} or
a separate \verb|\bibliographystyle| command on top of this setup.
